\begin{note}
	Также, как и для функций одной переменной, из диффренцируемости $f,\ g$ в точке $\v x_0$ следует дифференцируемость в точке $\v x_0$ следующих штук:
	\[f\pm g,\ f\cdot g,\ \frac{f}{g},\ d(f\pm g)=df\pm dg,\ d(fg)=f\cdot dg+g\cdot df,\ d(\frac fg)=\frac{g\cdot df-f\cdot dg}{g^2}\]
\end{note}
\begin{theorem}
	Если $f$ дифференцируема в точке $\v y\in\R^m$, а функции $g_j(\v x)$ дифференцируемы в точке $\v x_0\in\R^n$, причём $g_j(\v x_0)=y_{j,0},\ \v y_0=(y_{1,0},\dots,y_{m,0}),\ j=1,\dots,m$, то сложная функция $h(\v x)=f(g_1(x),\dots,g_m(\v x))$ дифференцируема в точке $\v x_0$, причём:
	\[\grad h(x)=\grad f(\v y_0)\cdot(\grad g_1(\v x_0),\dots,\grad g_m(\v x_0))^T\]
\end{theorem}
\begin{proof}
	$\tr y_j:=g_j(\v x_0+\tr\v x)-g_j(\v x_0),\ j=1,\dots,m$\\
	$\tr h(\v x_0)=f(g_1(\v x_0+\tr\v x),\dots,g_m(\v x_0+\tr\v x))-f(g_1(\v x_0),\dots,g_m(\v x_0))=\tr f(\v x_0)$\\
	Уточним, так как $\tr\v y$ может быть равен $\v0$, то доопределим $\sigma(0)=0$.
	\begin{multline*}
		\tr f(\v x_0)=(\grad f(\v y_0),\tr\v y)+\sigma(|\tr\v y|)=\\
		=\grad f(\v y_0)\cdot(\grad (g_1(\v x_0),\tr\v x)+\sigma(\tr\v x),\dots,\grad(g_m(\v x_0)+\sigma(\tr\v x),\tr\v x))^T+\sigma(|\tr\v y|)=\\
		=\underbrace{(\grad f(\v y_0))\cdot(\grad(g_1(\v x_0),\tr\v x),\dots,\grad(g_n(\v x_0),\tr\v x))^T}_{\grad(h(\v x_0),\tr\v x)}+\\
		+\underbrace{\grad(f(\v y_0))\cdot(\sigma_1(|\tr\v x|),\dots,\sigma_n(|\tr\v x|))^T+\sigma(|\tr\v y|)}_{\sigma(|\tr\v x|)}
	\end{multline*}
	\[\frac{\tr y_j}{|\tr\v x|}\le\frac{(\grad g_j(\v x_0),\tr\v x)}{|\tr\v x|}+\grad{|\sigma_j(|\tr\v x|)|}{|\tr\v x|},\ \frac{\sigma(|\tr\v y|)}{|\tr\v y|}\cdot\frac{\tr\v y}{\tr\v x}\ra0\]
	Получили представление $\grad h(x)$ в требуемом виде.
\end{proof}
\begin{corollary}[Инвариантность формы первого дифференциала]
	Формула для дифференциала функции $f(\v x):\ df=\sum_{i=1}^n\frac{\vd f}{\vd x_j}dx_j$ справедлива как в случае независимых переменных, так и в случае зависимых.
\end{corollary}
\begin{proof}
	\[f(\v x)=f(g_1(\v t),\dots,g_n(\v t))=:h(\v t)\]
	\[df=dh=\sum_{k=1}^m\frac{\vd h}{\vd t_k}dt_k=\sum_{k=1}^m\sum_{j=1}^n\frac{\vd f}{\vd x_j}\cdot\frac{\vd g}{\vd t_k}dt_k=\sum_{j=1}^n\frac{\vd f}{\vd x_j}(\sum_{k=1}^m\frac{\vd g_j}{\vd t_k}dt_k)=\sum_{j=1}^n\frac{\vd f}{\vd x_j}dg_j=\sum_{j=1}^n\frac{\vd f}{\vd x_j}dx_j\]
\end{proof}

\section{Геометрический смысл градиента и дифференцируемости}
\subsection{Параметрически заданная поверхность}
\begin{definition}
	\textit{Область} - открытое связное множество.
\end{definition}
\begin{definition}
	\textit{Замкнутая область} - замыкание области.
\end{definition}
\begin{definition}
	\textit{Параметрически заданной поверхностью} называется множество \\
	точек пространства $\R^3$, заданное с помощью непрерывных функций вида \\
	$x=x(u,v)$, $y=y(u,v)$, $z=z(u,v)$ на замкнутой области $D\subset\R^2$.
\end{definition}
\begin{definition}
	Плоскость, проходящая через точку $N_0$ параметрически заданной поверхности называется \textit{касательной} к этой поверхности, если угол между ней и любой секущей проходящей через точки $N_0$ и $N$ поверхности $\ra0$ при $N\ra N_0$. 
\end{definition}
\begin{definition}\label{def1}
	\textit{График функции}: $z=f(x,y)$ на $\v\Omega$, $\{(x,y,z):\ (x,y)\in\v\Omega,\ z=f(x,y)\}$
\end{definition}
\begin{note}
	К определению{~\eqref{def1}}: для точек на этой поверхности есть биекция с плоскостью $Ox$ и $Oy$.
\end{note}
\begin{theorem}
	Если $f(x,y)$ дифференцируема в точке $(x_0,y_0)$, то её график имеет касательную плоскость в точке $(x_0,y_0,z_0)$, $z_0=f(x_0,y_0)$, заданную уравнением:
	\[z-z_0=\frac{\vd f}{\vd x}(x_0,y_0)(x-x_0)+\frac{\vd f}{\vd y}(x_0,y_0)(y-y_0)\]
\end{theorem}
\begin{proof}
	Заметим, что заданная плоскость проходит через точку $N_0(x_0,y_0,z_0)$.\\
	Нормаль к ней имеет координаты $(\frac{\vd f}{\vd x}(x_0,y_0),\frac{\vd f}{\vd y}(x_0,y_0),-1)=\v n$\\
	$\v{N_0N}=(x-x_0,y-y_0,f(x,y)-f(x_0,y_0))$\\
	Хотим доказать:
	\[\frac{(\v n,\v{N_0N})}{|\v n\cdot|\v{N_0N}|}\ra0 \text{ при }(x-x_0,y-y_0)\ra0\]
	\[
	\left|\frac{\frac{\vd f}{\vd x}(x_0,y_0)(x-x_0)+\frac{\vd f}{\vd y}(x_0,y_0)(y-y_0)-(f(x,y)-f(x_0,y_0))}{|\v n|\cdot\sqrt{(x-x_0)^2+(y-y_0)^2+(f(x,y)-f(x_0,y_0))^2}}\right|\le\frac{\sigma(|\tr\v x|)}{|\tr\v x|}
	\]
	Получили требуемое.
\end{proof}
\subsection{Производная по направлению}
\begin{definition}
	\textit{Производной по направлению} $\v l\in\R^n\setminus\{\v0\}$ функции $f$ в точке $x_0\in\R^n$ называется предел (если он $\exists$):
	\[\frac{\vd f}{\vd l}(\v x_0)=\lim_{t\ra+0}\frac{f(\v x_0+t\cdot\v l)-f(\v x_0)}{t}\] 
\end{definition}
\begin{theorem}
	Если $f$ дифференцируема в точке $\v x_0\in\R^n$, то она имеет производную в этой точке по всем направлениям $\v l\in\R^n\setminus\{\v0\}$, причём:
	\[\frac{\vd f}{\vd\v l}(\v x_0)=(\grad f(\v x_0),\v l)\]
\end{theorem}
\begin{proof}
	\[\frac{f(\v x_0+t\cdot\v l)-f(\v x_0)}{t}=\frac{\grad f(\v x_0)+\sigma(|t\cdot\v l|)}{t}=(\grad f(\v x_0),\v l)+\sigma(1)\]
\end{proof}